% Part: turing-machines
% Chapter: machines-computations
% Section: representing-tms

\documentclass[../../../include/open-logic-section]{subfiles}

\begin{document}

\olfileid{tur}{mac}{rep}
\olsection{Nggambarake Mesin Turing}

\begin{explain}
Mesin Turing bisa digambarake sacara visual nganggo
\emph{diagram kahanan}. Diagram iki dumadi saka
kothak kahanan kang disambung nganggo panah. Mesthi wae,
saben kothak kahanan nggambarake kahanan mesin.
Saben panah nggambarake instruksi kang bisa ditindakake
saka kahanan iku, kanthi rincian instruksi ditulis ing
ndhuwur utawa ngisor panah kang cocog. Gatekna mesin
ing ngisor iki, kang mung nduwé rong kahanan internal,
$q_0$ lan $q_1$, lan siji instruksi:
\[
\begin{tikzpicture}[->,>=stealth',shorten >=1pt,auto,node distance=2.8cm,
                    semithick]
  \tikzstyle{every state}=[fill=none,draw=black,text=black]

  \node[initial,state]   (A)              {$q_0$};
  \node[state]   (B) [right of=A] {$q_1$};

  \path (A) edge  node {\TMtrans{\TMblank}{\TMstroke}{\TMright}} (B);
\end{tikzpicture}
\]
Elinga yen mesin Turing nduwé sirah maca-nulis lan
pita kang wis ditulisi input. Instruksi iku bisa
diwaca minangka \emph{yen maca~$\TMblank$ ing
kahanan $q_0$, tulisen~$\TMstroke$, obaha menyang
tengen, lan alihana menyang kahanan $q_1$}.
Iki ekuivalen karo fungsi transisi kang ngirim
$\tuple{q_0, \TMblank}$ menyang $\tuple{q_1, \TMstroke,
\TMright}$.
\end{explain}

\begin{ex}
\emph{Mesin Genep}: Mesin Turing ing ngisor iki
mandheg yen lan mung yen cacah $\TMstroke$ ing
pita genep (kanthi anggepan kabeh $\TMstroke$
ana sadurunge $\TMblank$ kapisan ing pita). Cathetan panyunting: conto iki miwiti sirah ing kothak input kapisan ing sisih tengen panandha, manut definisi formal konfigurasi wiwitan. SOL6-F049 uga mbenerake posisi wiwitan ing pambuka supaya konvensine padha. Kanggo input kosong sirah maca blank kapisan ing kahanan nol.
\[
\begin{tikzpicture}[->,>=stealth',shorten >=1pt,auto,node distance=2.8cm,
                    semithick]
  \tikzstyle{every state}=[fill=none,draw=black,text=black]

  \node[initial,state]         (A)              {$q_0$};
  \node[state]         (B) [right of=A] {$q_1$};

  \path (A) edge [bend left] node {\TMtrans{\TMstroke}{\TMstroke}{\TMright}} (B)
        (B) edge [loop above] node {\TMtrans{\TMblank}{\TMblank}{\TMright}} (B)
            edge [bend left] node {\TMtrans{\TMstroke}{\TMstroke}{\TMright}} (A);
\end{tikzpicture}
\]

Diagram kahanan iku cocog karo fungsi transisi
ing ngisor iki:
\begin{align*}
\delta(q_0, \TMstroke) & = \tuple{q_1, \TMstroke, \TMright},\\
\delta(q_1, \TMstroke) & = \tuple{q_0, \TMstroke, \TMright},\\
\delta(q_1, \TMblank)  & = \tuple{q_1, \TMblank, \TMright}
\end{align*}
\end{ex}

\begin{explain}
Mesin ing ndhuwur mung mandheg yen inpute
cacah guratan kang genep. Yen ora,
mesin iku (miturut teori) terus lumaku
tanpa wates. Kanggo mesin lan input
sembarang, kita bisa nelusuri
\emph{konfigurasi} mesin kanggo nemtokake
asile yen mesin mandheg. Cathetan koreksi sumber OLPL-770: panelusuran menehi output yen mesin mandheg, nanging ora mesthi nemtokake output utawa mandhege mesin kanggo input sembarang. Kita bakal menehi definisi formal
konfigurasi ing bagean sabanjure. Saiki,
konfigurasi bisa dianggep sacara intuisi
minangka rerangken diagram kang nuduhake
kahanan mesin ing saben wektu nalika lumaku.
Konfigurasi nuduhake isi pita, kahanan mesin,
lan panggonan sirah maca-nulis.

Ayo kita telusuri konfigurasi mesin genep
yen diwiwiti nganggo input patang
$\TMstroke$. Ing kahanan iki, kita ngarepake
mesin bakal mandheg. Banjur kita bakal
nglakokake mesin nganggo input telung
$\TMstroke$, kang bakal njalari mesin
lumaku tanpa entek.

Mesin diwiwiti ing kahanan~$q_0$ lan maca
$\TMstroke$ paling kiwa. Kita bisa
nggambarake kahanan wiwitan mesin kaya mangkene:
\[
\TMendtape \TMstroke_0 \TMstroke \TMstroke \TMstroke \TMblank \ldots
\]
% OLPL-145: Sumber beku nyebut "state one" ing katrangan konfigurasi wiwitan, nanging rumus lan diagram nuduhake q_0.
Konfigurasi ing ndhuwur gampang dimangerteni.
Cathetan koreksi sumber OLPL-145: kahanan wiwitan nol, dudu siji kaya ing ukara sumber. Kaya kang katon, mesin diwiwiti ing kahanan
nol lan maca~$\TMstroke$ paling kiwa.
Iki digambarake dening subskrip jeneng
kahanan ing $\TMstroke$ kapisan.
Instruksi kang ditrapake ing kene yaiku
$\delta(q_0, \TMstroke) = \tuple{q_1, \TMstroke, \TMright}$,
mula mesin obah menyang tengen ing pita lan
ngalih menyang kahanan~$q_1$.
\[
\TMendtape \TMstroke \TMstroke_1 \TMstroke \TMstroke \TMblank \ldots
\]
Amarga mesin saiki ana ing kahanan~$q_1$
lan maca~$\TMstroke$, kita kudu
``ngetutake'' instruksi $\delta(q_1, \TMstroke) = \tuple{q_0,
  \TMstroke, \TMright}$. Iki ngasilake konfigurasi
\[
\TMendtape \TMstroke \TMstroke \TMstroke_0 \TMstroke \TMblank \ldots
\]
Nalika mesin nerusake, aturan-aturan ditrapake
maneh kanthi urutan kang padha, ngasilake
rong konfigurasi ing ngisor iki:
\[
\TMendtape \TMstroke \TMstroke \TMstroke \TMstroke_1 \TMblank \ldots
\]
\[
\TMendtape \TMstroke \TMstroke \TMstroke \TMstroke \TMblank_0 \ldots
\]
Mesin saiki ana ing kahanan~$q_0$ lan maca
$\TMblank$. Adhedhasar diagram transisi,
gampang dideleng yen ora ana instruksi
kang bisa ditindakake, mula mesin mandheg.
Iki tegese input ditampa.

Saiki upamane mesin diwiwiti nganggo input
telung $\TMstroke$. Sawetara konfigurasi
wiwitan padha karo sadurunge, amarga
instruksi kang padha ditindakake, mung
input ing pita kang rada beda:
\[
\TMendtape \TMstroke_0 \TMstroke \TMstroke \TMblank \ldots
\]
\[
\TMendtape \TMstroke \TMstroke_1 \TMstroke \TMblank \ldots
\]
\[
\TMendtape \TMstroke \TMstroke \TMstroke_0 \TMblank \ldots
\]
\[
\TMendtape \TMstroke \TMstroke \TMstroke \TMblank_1 \ldots
\]
Saiki mesin wis ngliwati kabeh $\TMstroke$
lan maca~$\TMblank$ ing kahanan~$q_1$.
Kaya kang katon ing diagram, ana instruksi
$\delta(q_1, \TMblank) =\tuple{q_1,
\TMblank, \TMright}$. Amarga pita diisi
$\TMblank$ tanpa wates menyang tengen,
mesin bakal terus nglakokake instruksi iki
\emph{salawase}, tetep ana ing kahanan~$q_1$
lan obah saya adoh menyang tengen. Mesin
ora bakal mandheg lan ora nampa input iku.
\end{explain}

\begin{explain}
Perlu digatekake yen ora kabeh mesin bakal
mandheg. Yen mandheg tegese mesin ora
nduwé instruksi maneh kang bisa ditindakake,
kita bisa nggawe mesin kang ora bakal
mandheg kanthi mesthekake ana panah metu
kanggo saben simbol ing saben kahanan.
Mesin genep bisa diowahi supaya lumaku
tanpa wates kanthi nambah instruksi
kanggo maca $\TMblank$ ing~$q_0$.
\end{explain}

\begin{ex}
\[
\begin{tikzpicture}[->,>=stealth',shorten >=1pt,auto,node distance=2.8cm,
                    semithick]
  \tikzstyle{every state}=[fill=none,draw=black,text=black]

  \node[initial,state] (A)              {$q_0$};
  \node[state]         (B) [right of=A] {$q_1$};

  \path
  (A) edge [bend left] node {\TMtrans{\TMstroke}{\TMstroke}{\TMright}} (B)
      edge [loop above] node {\TMtrans{\TMblank}{\TMblank}{\TMright}} (A)
  (B) edge [loop above] node {\TMtrans{\TMblank}{\TMblank}{\TMright}} (B)
      edge [bend left] node {\TMtrans{\TMstroke}{\TMstroke}{\TMright}} (A);
\end{tikzpicture}
\]
\end{ex}

\begin{explain}
Tabel mesin minangka cara liyane kanggo
nggambarake mesin Turing. Tabel mesin
nuduhake alfabet pita ing sumbu~$x$
lan himpunan kahanan mesin ing sumbu~$y$.
Ing tabel, ing pepanggihan saben kahanan
lan simbol, ditulis bagean instruksi liyane
---kahanan anyar, simbol anyar, lan arah
obah. Tabel mesin nggampangake nemtokake
ing kahanan lan simbol apa mesin bakal
mandheg. Saben papan kosong ing tabel
minangka titik kang bisa njalari mesin
mandheg. Beda karo diagram kahanan lan
dhaptar instruksi, kang titik-titik
mandhege mesin ora mesthi langsung cetha,
titik-titik mandheg ing tabel mesin bisa
cepet ditemokake saka papan-papan kosong.
\end{explain}

\begin{ex}
Tabel mesin kanggo mesin genep yaiku:
\[
\centering
\begin{tabular}{lllll}
\cline{2-4}
\multicolumn{1}{l|}{}      & \multicolumn{1}{c|}{$\TMblank$}                
& \multicolumn{1}{c|}{$\TMstroke$}     & \multicolumn{1}{c|}{$\TMendtape$}   &  \\ \cline{1-4}
\multicolumn{1}{|l|}{$q_0$} & \multicolumn{1}{l|}{}                          
& \multicolumn{1}{l|}{\TMtrans{\TMstroke}{q_1}{\TMright}} & 
\multicolumn{1}{l|}{\phantom{\TMtrans{\TMstroke}{q_1}{\TMright}}} &  \\ \cline{1-4}
\multicolumn{1}{|l|}{$q_1$} & \multicolumn{1}{l|}{\TMtrans{\TMblank}{q_1}{\TMright}} 
& \multicolumn{1}{l|}{\TMtrans{\TMstroke}{q_0}{\TMright}} & 
\multicolumn{1}{l|}{\phantom{\TMtrans{\TMstroke}{q_1}{\TMright}}} &  \\ \cline{1-4}
\end{tabular}
\]
Kaya kang katon, mesin mandheg nalika
maca~$\TMblank$ ing kahanan~$q_0$.
\end{ex}

\begin{explain}
Nganti saiki kita mung nggatekake mesin kang
maca lan nampa input. Nanging mesin Turing
bisa maca lan nulis. Salah siji tuladhane
(sanajan ana akeh banget) yaiku
\emph{mesin pangganda}. Mesin pangganda
kang diwiwiti nganggo blok~$n$
$\TMstroke$ ing pita ngasilake blok~$2n$
$\TMstroke$.
\end{explain}

\begin{ex}
  \ollabel{ex:doubler} Sadurunge mbangun
  mesin pangganda, kita kudu nemtokake
  \emph{strategi} kanggo ngrampungake masalah.
  Amarga mesin (miturut rumusan kita) ora
  bisa ngelingi cacah $\TMstroke$ kang wis
  diwaca, kita butuh cara kanggo ngetutake
  kabeh $\TMstroke$ ing pita. Salah siji
  carane yaiku misahake output saka input
  nganggo~$\TMblank$. Mesin banjur bisa
  mbusak $\TMstroke$ kapisan saka input,
  ngliwati input sisane, ninggalake
  $\TMblank$, banjur nulis rong $\TMstroke$
  anyar. Sabanjure mesin bali lan nemokake
  $\TMstroke$ kapindho ing input, banjur
  nggandakake guratan iku uga. Kanggo saben
  siji $\TMstroke$ ing input, mesin nulis
  rong $\TMstroke$ ing output. Kanthi
  mbusak input nalika mesin lumaku, kita
  bisa mesthekake ora ana $\TMstroke$
  kang kliwatan utawa digandakake kaping
  pindho. Nalika kabeh input wis kabusak,
  ing pita kari $2n$ $\TMstroke$.
  Diagram kahanan mesin Turing kang diasilake
  digambarake ing \olref{fig:doubler}.
  \begin{figure}
\[
\begin{tikzpicture}[->,>=stealth',shorten >=1pt,auto,node distance=2.8cm,
                    semithick]
  \tikzstyle{every state}=[fill=none,draw=black,text=black]

  \node[initial,state] (1)              {$q_0$};
  \node[state]         (2) [right of=1] {$q_1$};
  \node[state]         (3) [right of=2] {$q_2$};
  \node[state]         (4) [below of=3] {$q_3$};
  \node[state]         (5) [left of=4]  {$q_4$};
  \node[state]         (6) [left of=5]  {$q_5$};

  \path (1) edge node {\TMtrans{\TMstroke}{\TMblank}{\TMright}} (2)
    (2) edge [loop above] node {\TMtrans{\TMstroke}{\TMstroke}{\TMright}} (2)
      edge node {\TMtrans{\TMblank}{\TMblank}{\TMright}} (3)
    (3) edge [loop above] node {\TMtrans{\TMstroke}{\TMstroke}{\TMright}} (3)
        edge node {\TMtrans{\TMblank}{\TMstroke}{\TMright}} (4)
    (4) edge [loop below] node {\TMtrans{\TMblank}{\TMstroke}{\TMleft}} (4)
        edge node {\TMtrans{\TMstroke}{\TMstroke}{\TMleft}} (5)
    (5) edge [loop below]  node {\TMtrans{\TMstroke}{\TMstroke}{\TMleft}} (5)
        edge              node {\TMtrans{\TMblank}{\TMblank}{\TMleft}} (6)
    (6) edge [loop below] node {\TMtrans{\TMstroke}{\TMstroke}{\TMleft}} (6)
        edge              node {\TMtrans{\TMblank}{\TMblank}{\TMright}} (1);
\end{tikzpicture}
\]
\caption{Mesin pangganda}
\ollabel{fig:doubler}
\end{figure}
\end{ex}

\begin{prob}
Pilihen input sembarang lan telusurana konfigurasi
mesin pangganda ing \olref[tur][mac][rep]{ex:doubler}.
\end{prob}

\begin{prob}
Rancangna mesin Turing kanthi alfabet
$\{\TMendtape,\TMblank, A, B\}$ kang nampa,
yaiku mandheg nalika nampa, rerangken aksara
$A$ lan $B$ sembarang kanthi cacah $A$
padha karo cacah $B$ \emph{lan} kabeh $A$
ana sadurunge kabeh $B$; lan kang nolak,
yaiku ora mandheg nalika nampa, rerangken
aksara kang cacah $A$ ora padha karo cacah~$B$
utawa ora kabeh $A$ ana sadurunge kabeh~$B$.
(Tuladhane, mesin kudu nampa $AABB$ lan
$AAABBB$, nanging nolak $AAB$ lan $AABBAABB$.)
\end{prob}

\begin{prob}
Rancangna mesin Turing kanthi alfabet
$\{\TMendtape,\TMblank, A, B\}$ kang nampa
rerangken aksara $\alpha$ sembarang saka $A$
lan $B$ minangka input lan nggandakake
rerangken iku supaya ngasilake $\alpha\alpha$.
(Tuladhane, input $ABBA$ kudu ngasilake
output $ABBAABBA$.)
\end{prob}

\begin{prob}
\emph{Miturut urutan alfabet?:} Rancangna
mesin Turing kanthi alfabet
$\{\TMendtape,\TMblank, A, B\}$ kang nalika
diwenehi input rerangken winates saka $A$
lan $B$ mriksa apa kabeh $A$ ana ing
sisih kiwa kabeh $B$. Mesin kudu ninggalake
rerangken input ing pita, lan mandheg yen
rerangkene ``miturut urutan alfabet'',
utawa lumaku ing daur tanpa wates yen ora.
\end{prob}

\begin{prob}
\emph{Pangurut alfabet:} Rancangna mesin
Turing kanthi alfabet
$\{\TMendtape,\TMblank, A, B\}$ kang nampa
rerangken winates saka $A$ lan $B$ minangka
input banjur nata maneh supaya kabeh $A$
ana ing sisih kiwa kabeh~$B$. (Tuladhane,
rerangken $BABAA$ kudu dadi rerangken $AAABB$,
lan rerangken $ABBABB$ kudu dadi rerangken
$AABBBB$.)
\end{prob}

\end{document}
